\subsection{Multi-Tier Iterative Model Discovery Architecture}

To provide users with the most appropriate AI models for their specific technical requirements, the framework implements a Two-Tier Model Discovery architecture. This system manages the transition from curated pre-defined models to broader, automated web-based retrieval.

\subsubsection{Tier 1: Task-Based Curated Selection}

The first discovery stage relies on a task-based query extraction. When a user describes a project (e.g., "I want to perform real-time image classification"), the system performs two simultaneous actions:
\begin{enumerate}
    \item \textbf{LLM Extraction}: Task identification (e.g., \texttt{image\_classification}) and hardware profiling.
    \item \textbf{Zoo Lookup}: The system queries a local registry of optimized models sourced from the \textit{STM32 AI Model Zoo}.
\end{enumerate}

The user is presented with a list of curated models (e.g., MobileNetV2, SqueezeNet) that are known to be compatible with the STM32 ecosystem.

\subsubsection{Tier 2: Iterative Search Loop}

If the user rejects the curated list or requires a niche architecture not present in the local registry, the system transitions into an iterative search loop. This loop consists of up to three attempts to find a suitable \texttt{.h5} model file through two search backends:

\begin{itemize}
    \item \textbf{GitHub Search (Primary)}: Hybrid search in Model Zoo repositories for \texttt{.h5} files.
    \item \textbf{Google Search Fallback}: If the GitHub search fails to yield a direct model link, the system broadened the search using Google Search tools to locate community-contributed models or technical documentation.
\end{itemize}

\subsubsection{Strategic Fallback and Failsafe}

The discovery process is governed by a strict failsafe logic to ensure the workflow can always proceed to the optimization stage. If the iterative search exhausts the three attempts without user acceptance, the system automatically defaults to a hierarchical fallback:
\begin{enumerate}
    \item \textbf{Task-Based Fallback}: The primary optimized model for the identified task is selected from the local registry.
    \item \textbf{System Default}: If no task can be clearly identified, the system utilizes a global baseline model (e.g., standard MobileNetV2) to ensure continuity.
\end{enumerate}

This robust discovery architecture balances user preference with automated reliability, ensuring that the development pipeline is never stalled due to a lack of initial model weights.
